% !TEX root = ./main.tex

% ============================================================
% 五人团队信息：在这里填写姓名、学号和导师
% 英文姓名请使用学籍系统中的正式拼写
% ============================================================

% 第一位成员
\newcommand{\TeamMemberOneZh}{杨皓然}
\newcommand{\TeamMemberOneEn}{Yang Haoran}
\newcommand{\TeamMemberOneID}{520XXXXXXXX}

% 第二位成员
\newcommand{\TeamMemberTwoZh}{吴律}
\newcommand{\TeamMemberTwoEn}{Wu Lu}
\newcommand{\TeamMemberTwoID}{520XXXXXXXX}

% 第三位成员
\newcommand{\TeamMemberThreeZh}{吴毅昕}
\newcommand{\TeamMemberThreeEn}{Wu Yixin}
\newcommand{\TeamMemberThreeID}{520XXXXXXXX}

% 第四位成员
\newcommand{\TeamMemberFourZh}{程天纵}
\newcommand{\TeamMemberFourEn}{Cheng Tianzong}
\newcommand{\TeamMemberFourID}{520XXXXXXXX}

% 第五位成员
\newcommand{\TeamMemberFiveZh}{顾益铭}
\newcommand{\TeamMemberFiveEn}{Gu Yiming}
\newcommand{\TeamMemberFiveID}{520XXXXXXXX}

% 其他封面信息
\newcommand{\ThesisSupervisorZh}{李四}
\newcommand{\ThesisSupervisorEn}{Li Si}

\newcommand{\ThesisDepartmentZh}{浦江国际学院}
\newcommand{\ThesisDepartmentEn}{Global College}

\newcommand{\ThesisMajorZh}{工业工程}
\newcommand{\ThesisMajorEn}{Industrial Engineering}

\newcommand{\ThesisDegreeZh}{学士}

% 用于论文元数据的姓名和学号列表
\newcommand{\TeamNamesZhInline}{%
  \TeamMemberOneZh、\TeamMemberTwoZh、\TeamMemberThreeZh、%
  \TeamMemberFourZh、\TeamMemberFiveZh%
}

\newcommand{\TeamNamesEnInline}{%
  \TeamMemberOneEn, \TeamMemberTwoEn, \TeamMemberThreeEn, %
  \TeamMemberFourEn, and \TeamMemberFiveEn%
}

\newcommand{\TeamIDsInline}{%
  \TeamMemberOneID, \TeamMemberTwoID, \TeamMemberThreeID, %
  \TeamMemberFourID, \TeamMemberFiveID%
}

% 中文封面的五人姓名和学号表
\newcommand{\TeamMembersZh}{%
  \begingroup
  \zihao{4}%
  \renewcommand{\arraystretch}{1.08}%
  \begin{tabular}[t]{@{}l@{\hspace{3.5em}}l@{}}
    \heiti 姓名 & \heiti 学号 \\
    \TeamMemberOneZh   & \TeamMemberOneID   \\
    \TeamMemberTwoZh   & \TeamMemberTwoID   \\
    \TeamMemberThreeZh & \TeamMemberThreeID \\
    \TeamMemberFourZh  & \TeamMemberFourID  \\
    \TeamMemberFiveZh  & \TeamMemberFiveID
  \end{tabular}%
  \endgroup
}

% 英文封面的五人姓名和学号表
\newcommand{\TeamMembersEn}{%
  \begingroup
  \zihao{4}%
  \renewcommand{\arraystretch}{1.08}%
  \begin{tabular}[t]{@{}l@{\hspace{3.5em}}l@{}}
    \textbf{Name} & \textbf{Student ID} \\
    \TeamMemberOneEn   & \TeamMemberOneID   \\
    \TeamMemberTwoEn   & \TeamMemberTwoID   \\
    \TeamMemberThreeEn & \TeamMemberThreeID \\
    \TeamMemberFourEn  & \TeamMemberFourID  \\
    \TeamMemberFiveEn  & \TeamMemberFiveID
  \end{tabular}%
  \endgroup
}

% 中文封面的导师、学院、专业和学位
\newcommand{\CoverDetailsZh}{%
  \begingroup
  \zihao{4}%
  \renewcommand{\arraystretch}{1.08}%
  \begin{tabular}[t]{@{}r@{\hspace{0.5em}：\hspace{0.8em}}l@{}}
    \heiti 导\qquad 师 & \ThesisSupervisorZh \\
    \heiti 学\qquad 院 & \ThesisDepartmentZh \\
    \heiti 专业名称 & \ThesisMajorZh \\
    \heiti 申请学位层次 & \ThesisDegreeZh
  \end{tabular}%
  \endgroup
}

% 英文封面的导师
\newcommand{\CoverSupervisorEn}{%
  \begingroup
  \zihao{4}%
  \begin{tabular}[t]{@{}r@{\hspace{0.5em}:\hspace{0.8em}}l@{}}
    \textbf{Supervisor} & \ThesisSupervisorEn
  \end{tabular}%
  \endgroup
}

\sjtusetup{
  %
  %******************************
  % 注意：
  %   1. 配置里面不要出现空行
  %   2. 不需要的配置信息可以删除
  %******************************
  %
  info = {%
    % 论文题目
    zh / title = {
      面向 3C 工业生产线的双臂机器人视觉-语言-动作模型
    },
    en / title = {
      A Vision-Language-Action Model for Dual-Arm Robots
      in 3C Industrial Production Lines
    },
    % 关键词
    zh / keywords = {
      双臂机器人，视觉-语言-动作模型，工业自动化
    },
    en / keywords = {
      dual-arm robot, vision-language-action model,
      industrial automation
    },
    % 作者
    zh / author = {\TeamNamesZhInline},
    en / author = {\TeamNamesEnInline},
    % 导师
    zh / supervisor = {\ThesisSupervisorZh},
    en / supervisor = {\ThesisSupervisorEn},
    % 五个学号
    id = {\TeamIDsInline},
    % 学位
    zh / degree = {\ThesisDegreeZh},
    en / degree = {Bachelor of Engineering},
    % 专业
    zh / major = {\ThesisMajorZh},
    en / major = {\ThesisMajorEn},
    % 学院
    zh / department = {\ThesisDepartmentZh},
    en / department = {\ThesisDepartmentEn},
    % 日期
    date = {2026-08},
    zh / display-date = {2026 年 8 月},
  },
  style = {%
    % title-logo-color = black,
  },
  name = {
    digest = {},
  },
}

% ============================================================
% 将原来的单人封面信息区改为五人表格
% 只修改 setup.tex，不需要修改 cls 或 def 文件
% ============================================================

\ExplSyntaxOn

% 中文封面
\EditInstance { sjtu } { title / zh / info }
  {
    format = \zihao { 4 } \fixedlineskip { 31.2 bp },

    content =
      {
        \TeamMembersZh

        \par
        \vspace { 7.8 bp }

        \CoverDetailsZh
      },

    bottom-skip = 31.2 bp
  }

% 英文封面
\EditInstance { sjtu } { title / en / info }
  {
    format = \zihao { 3 } \fixedlineskip { 31.2 bp },

    content =
      {
        \TeamMembersEn

        \par
        \vspace { 7.8 bp }

        \CoverSupervisorEn
      },

    bottom-skip = \c_zero_dim plus 3 fill
  }

\ExplSyntaxOff

% 使用 BibLaTeX 处理参考文献
%   biblatex-gb7714-2015 常用选项
%     gbnamefmt=lowercase     姓名大小写由输入信息确定
%     gbpub=false             禁用出版信息缺失处理
\usepackage[backend=biber,style=gb7714-2015]{biblatex}
% 文献表字体
% \renewcommand{\bibfont}{\zihao{5}\fixedlineskip{15.6bp}}
% 文献表条目间的间距
\setlength{\bibitemsep}{0pt}
% 导入参考文献数据库
\addbibresource{refs.bib}

% 脚注格式
\usepackage[perpage,bottom,hang]{footmisc}

% 定义图片文件目录与扩展名
\graphicspath{{figures/}}
\DeclareGraphicsExtensions{.pdf,.eps,.png,.jpg,.jpeg}

% 确定浮动对象的位置，可以使用 [H]，强制将浮动对象放到这里（可能效果很差）
% \usepackage{float}

% 固定宽度的表格
% \usepackage{tabularx}

% 使用三线表：toprule，midrule，bottomrule。
\usepackage{booktabs}

% 表格中支持跨行
\usepackage{multirow}

% 表格中数字按小数点对齐
\usepackage{dcolumn}
\newcolumntype{d}[1]{D{.}{.}{#1}}

% 使用长表格
\usepackage{longtable}

% 附带脚注的表格
\usepackage{threeparttable}

% 附带脚注的长表格
\usepackage{threeparttablex}

% 算法环境宏包
\usepackage[ruled,vlined,linesnumbered]{algorithm2e}
% \usepackage{algorithm, algorithmicx, algpseudocode}

% 代码环境宏包
\usepackage{listings}
\lstdefinestyle{lstStyleCode}{%
  aboveskip         = \medskipamount,
  belowskip         = \medskipamount,
  basicstyle        = \ttfamily\zihao{6},
  commentstyle      = \slshape\color{black!60},
  stringstyle       = \color{green!40!black!100},
  keywordstyle      = \bfseries\color{blue!50!black},
  extendedchars     = false,
  upquote           = true,
  tabsize           = 2,
  showstringspaces  = false,
  xleftmargin       = 1em,
  xrightmargin      = 1em,
  breaklines        = false,
  framexleftmargin  = 1em,
  framexrightmargin = 1em,
  backgroundcolor   = \color{gray!10},
  columns           = flexible,
  keepspaces        = true,
  texcl             = true,
  mathescape        = true
}
\lstnewenvironment{codeblock}[1][]{%
  \lstset{style=lstStyleCode,#1}}{}

% 直立体数学符号
\providecommand{\dd}{\mathop{}\!\mathrm{d}}
\providecommand{\ee}{\mathrm{e}}
\providecommand{\ii}{\mathrm{i}}
\providecommand{\jj}{\mathrm{j}}

% 国际单位制宏包
\usepackage{siunitx}

% 定理环境宏包
\usepackage{ntheorem}
% \usepackage{amsthm}

% 绘图宏包
\usepackage{tikz}
\usetikzlibrary{arrows.meta, shapes.geometric}

% 一些文档中用到的 logo
\usepackage{hologo}
\providecommand{\XeTeX}{\hologo{XeTeX}}
\providecommand{\BibLaTeX}{\textsc{Bib}\LaTeX}

% 借用 ltxdoc 里面的几个命令方便写文档
\DeclareRobustCommand\cs[1]{\texttt{\char`\\#1}}
\providecommand\pkg[1]{{\sffamily#1}}

% hyperref 宏包在最后调用
\usepackage{hyperref}

% E-mail
\providecommand{\email}[1]{\href{mailto:#1}{\urlstyle{tt}\nolinkurl{#1}}}
